\section{SignalRule Module}
\label{sec:methodology:signal-rule}

The SignalRule module encapsulates the interlocking logic associated with a specific signal aspect condition, defining the circumstances under which a controlling signal permits another signal to display a given aspect. Each rule holds the necessary dependencies—such as track occupancy, point machine alignment, or protecting signal states—and provides methods to evaluate whether a requested aspect transition is allowed. By isolating these conditions into discrete rule objects, the system maintains modularity, enabling individual rules to be updated or extended without altering the overall engine.

When a validation query is initiated, the rule evaluates its stored conditions and determines whether the target signal may safely assume the requested aspect. To improve runtime efficiency, the module employs an internal cache that maps signal identifiers to their permitted aspects, reducing the overhead of repeated lookups during high-frequency operations. This approach ensures that real-time performance targets are maintained even in complex station environments with numerous signals and overlapping dependencies.

Integration with the InterlockingRuleEngine is seamless: rules are constructed from parsed JSON data and executed dynamically during validation cycles. Their outcomes feed into the broader validation framework, producing structured results that include both the decision (allowed or blocked) and the reasoning behind it. This explicit traceability supports debugging, audits, and compliance with safety standards, while the separation of rule logic from the core engine reinforces modularity and long-term maintainability. In practice, the SignalRule module serves as the building block of declarative interlocking logic, bridging abstract configuration data with executable safety enforcement in a clear and efficient manner.